In the last year alone, a surge of new benchmarks to measure \textit{compositional} understanding of vision-language models have permeated the machine learning ecosystem.
Given an image, these benchmarks probe a model's ability to identify its associated caption amongst a set of compositional distractors.
Surprisingly, we find significant biases in \textit{all} these benchmarks rendering them hackable. This hackability is so dire that blind models with no access to the image outperform state-of-the-art vision-language models.
To remedy this rampant vulnerability, we introduce \method, a new benchmark for vision-language compositionality evaluation.
We employ large language models, instead of rule-based templates used in previous benchmarks, to generate fluent and sensical hard negatives, and utilize an adversarial refinement mechanism to maximally reduce biases. We re-evaluate state-of-the-art models and recently proposed compositionality inducing strategies, and find that their improvements were hugely overestimated, suggesting that more innovation is needed in this important direction. We release \method and the code for evaluation at: \url{https://github.com/RAIVNLab/sugar-crepe}.